%--------------------
% Packages
% -------------------
\documentclass[11pt,a4paper]{article}
\usepackage[utf8x]{inputenc}
\usepackage[T1]{fontenc}
\usepackage{mathptmx} % Use Times Font

\usepackage[pdftex]{graphicx}
\usepackage[english]{babel} % Changed to English
\usepackage[pdftex,linkcolor=black,pdfborder={0 0 0}]{hyperref}
\usepackage{calc}
\usepackage{enumitem}
\usepackage{amsmath} % For math formatting
\usepackage{booktabs} % For professional tables
\usepackage{float} % For precise placement of figures/tables
\usepackage[utf8]{inputenc}
\usepackage{graphicx}
\usepackage{cite}
\usepackage{url}
\usepackage{booktabs}
\usepackage{hyperref}

\frenchspacing % No double spacing between sentences
\linespread{1.2} % Set linespace
% Enforcing strict 1cm margins as requested
\usepackage[a4paper, margin=1cm]{geometry} 

\usepackage[all]{nowidow}
\usepackage[protrusion=true,expansion=true]{microtype}

\hypersetup{ 	
pdfsubject = {Probabilistic Machine Learning},
pdftitle = {Link Prediction: Theory and Applications},
pdfauthor = {M.Tech Student}
}

\begin{document}

\begin{center}
    \LARGE\textbf{Bayesian Data Analysis}\vspace{0.5em}\\
    \large\textit{The ML Trifecta - Shalin Shah, Umesh BP, Ritik Agarwal} \\
    \large\textit{shalinshah@iisc.ac.in, umeshbp@iisc.ac.in, ritikagarwal@iisc.ac.in}
\end{center}
\vspace{1em}

\section{Introduction}
This project uses Bayesian methods to analyze the Job Salary Prediction dataset from Kaggle. It looks at how factors like experience, education, skills, and company details affect salary and remote work. We apply Bayesian Linear Regression, Logistic Classification, and Gaussian Mixture Clustering using PyMC. The models use simple priors, and we evaluate them using standard checks like convergence and model comparison.

\section{Data Description}
Dataset: \href{https://www.kaggle.com/datasets/nalisha/job-salary-prediction-dataset/data}{Job Salary Prediction Dataset}. The dataset contains 250,000 observations with 10 features and no missing values.

\begin{table}[htbp]
\centering
\begin{tabular}{lll}
\hline
\textbf{Feature} & \textbf{Type} & \textbf{Description} \\
\hline
job\_title & Categorical (12) & AI Engineer, Data Analyst, etc. \\
experience\_years & Numerical (0--20) & Years of professional experience \\
education\_level & Categorical (5) & High School to PhD \\
skills\_count & Numerical (1--19) & Number of skills \\
industry & Categorical (10) & Healthcare, Finance, etc. \\
company\_size & Categorical (5) & Startup to Enterprise \\
location & Categorical (10) & Country/region \\
remote\_work & Categorical (3) & No, Hybrid, Yes \\
certifications & Numerical (0--5) & Number of certifications \\
salary & Numerical & Target variable (31K--333K) \\
\hline
\end{tabular}
\end{table}

For computational efficiency, we used a random subset of 5,000 samples with standardized features and an 80/20 train-test split. To our knowledge, no prior Bayesian analysis exists for this dataset on Kaggle.

\section{Model Descriptions}

\subsection{Model 1: Bayesian Linear Regression (PyMC)}
A standard Bayesian linear regression on standardized features (experience, skills, certifications, education, company size) predicting standardized salary:

\[
\begin{aligned}
y      &\sim \text{Normal}(\alpha + X\beta, \sigma) \qquad\qquad & \alpha &\sim \text{Normal}(0, 10) \\
\beta  &\sim \text{Normal}(0, 5) \qquad\qquad & \sigma &\sim \text{HalfNormal}(5)
\end{aligned}
\]

\subsection{Model 2: Non-linear Bayesian Regression (PyMC)}
Extends Model 1 with interaction terms (experience $\times$ education, experience $\times$ company\_size) and a quadratic experience term ($\text{experience}^2$), yielding 8 predictors. Same prior structure.

\subsection{Model 3: Full-feature Bayesian Regression (PyMC)}
Extends the regression to include all categorical features (job\_title, industry, location) as one-hot encoded dummy variables alongside the 5 numerical predictors, yielding 34 features total. This model captures salary variation across job titles, industries, and locations:


\[
\begin{aligned}
y      &\sim \text{Normal}(\alpha + X\beta, \sigma) \qquad\qquad & \alpha &\sim \text{Normal}(0, 10) \\
\beta  &\sim \text{Normal}(0, 5) \qquad\qquad & \sigma &\sim \text{HalfNormal}(5)
\end{aligned}
\]

\noindent Training subsample: 4,000 observations for MCMC feasibility.

\vspace{0.5em}
\noindent \textbf{Implementation Details:} Categorical features (job\_title, industry, location) are one-hot encoded using \texttt{pandas.get\_dummies}, then concatenated with standardized numerical features. The resulting design matrix (34 predictors) is used in a PyMC model with weakly informative priors. Sampling is performed with NUTS (2 chains, 1000 draws, 500 tune). Model fit and coefficients are summarized using ArviZ.

\subsection{Model 4: Bayesian Logistic Regression (PyMC)}
Binary classification of remote work (Yes vs No, excluding Hybrid) using Bayesian logistic regression with the same 5 standardized features:
\vspace{-0.5em}
\begin{gather*}
P(\text{remote}=1 \mid X) = \text{sigmoid}(\alpha + X\beta) \\
\alpha \sim \text{Normal}(0, 5) \qquad \beta \sim \text{Normal}(0, 3)
\end{gather*}

\noindent \textbf{Implementation Details:} Hybrid rows are excluded and a binary target is created (remote=Yes:1, No:0). Standardized numerical features are used as predictors. The model is built in PyMC with a logistic link, weakly informative priors, and NUTS sampling (2 chains, 1000 draws, 500 tune). Posterior predictive checks and accuracy metrics are computed on the test set.

\subsection{Model 5: Bayesian Gaussian Mixture Clustering}
Bayesian Gaussian Mixture Model with Dirichlet Process prior (\texttt{sklearn}'s \texttt{BayesianGaussianMixture}) to discover natural groupings. Evaluated with $K=2$ to $7$ components; $K=4$ selected based on log-likelihood elbow.

\vspace{0.5em}
\noindent \textbf{Implementation Details:} Standardized numerical features are selected for clustering. The \texttt{BayesianGaussianMixture} from \texttt{sklearn} is fit for $K=2$ to $7$ components, using a Dirichlet process prior. Model selection is based on log-likelihood and effective number of components. Cluster assignments and profiles are analyzed for interpretation.

\section{Prior Choices and Justification}
We use weakly informative priors as the baseline. Since features are standardized, $\text{Normal}(0, 5)$ for coefficients allows moderate to large effect sizes without being overly constraining. $\text{HalfNormal}(5)$ for $\sigma$ is weakly informative for noise on the standardized scale. These priors regularize slightly but let the data dominate, which is appropriate for $N=4,000$.

\begin{table}[H]
\centering
\renewcommand{\arraystretch}{1.3}
\begin{tabular}{lll}
\toprule
\textbf{Prior Set} & \textbf{Coefficients (\(\beta\))} & \textbf{Noise (\(\sigma\))} \\
\midrule
\textbf{Weakly Informative (baseline)} & Normal(0, 5) & HalfNormal(5) \\
\textbf{Informative} & Normal(0, 1) & HalfNormal(1) \\
\textbf{Vague/Non-informative} & Normal(0, 100) & HalfCauchy(10) \\
\bottomrule
\end{tabular}
\end{table}

\section{Inference and Sampling}

All PyMC models were fit using the No-U-Turn Sampler (NUTS). For most models, we used 2 chains with 500 tuning iterations and 1,000 posterior draws per chain, with \texttt{target\_accept = 0.9} to improve robustness during sampling. Initialization was set to \texttt{jitter+adapt\_diag}. The linear regression model was initially run with 4 chains and 2,000 draws per chain.

\section{Convergence Diagnostics}

All models showed strong convergence behavior. The maximum R-hat values were approximately 1.000 across models, satisfying the common convergence guideline of being below 1.01. Effective sample sizes for the bulk of the posterior were also sufficiently large, and no divergent transitions were reported. In addition, trace plots indicated well-mixed chains across the fitted models.

\begin{table}[H]
\centering
\renewcommand{\arraystretch}{1.3}
\begin{tabular}{llll}
\toprule
\textbf{Model} & \textbf{R-hat (max)} & \textbf{ESS bulk (min)} & \textbf{Divergences} \\
\midrule
\textbf{Linear Regression} & 1.001 & 10,705 & 0 \\
\textbf{Non-linear Regression} & 1.000 & 1,746 & 0 \\
\textbf{Full-feature Regression} & 1.000 & 1,500 & 0 \\
\textbf{Logistic Classification} & 1.010 & 2,334 & 0 \\
\bottomrule
\end{tabular}
\caption{Convergence diagnostics for the fitted PyMC models.}
\label{tab:convergence_diagnostics}
\end{table}

\begin{figure}[H]
\centering
\includegraphics[width=0.85\textwidth]{linear-trace.png}
\caption{Trace plots or posterior diagnostic visualization.}
\label{fig:diagnostics}
\end{figure}

\section{Predictive Performance}

The full-feature Bayesian regression model achieved the strongest predictive performance among the regression models, with \(R^2 = 0.9552\). This result suggests that categorical and expanded feature information are the dominant drivers of salary variation in this dataset.

For classification, the Bayesian logistic regression model achieved an accuracy of 51.7\% for predicting remote work status. This is only marginally above the 50\% baseline, indicating that the available predictors---experience, skills, education, certifications, and company size---provide very limited information about whether a job is remote. The posterior coefficients were all close to zero, with wide highest density intervals spanning zero, which supports the conclusion that remote work is largely independent of these features in this dataset.

\section{Prior Sensitivity Analysis}

We compared three prior configurations on a 1,500-point subsample: weakly informative, informative, and vague priors. The posterior means remained highly stable across all prior choices, indicating that inference is robust to reasonable prior variation and is primarily driven by the observed data.

\begin{table}[H]
\centering
\renewcommand{\arraystretch}{1.3}
\begin{tabular}{lccc}
\toprule
\textbf{Parameter} & \textbf{Weakly Inf.} & \textbf{Informative} & \textbf{Vague} \\
\midrule
\(\beta_{\text{Experience}}\)     & 0.4401 & 0.4292 & 0.4291 \\
\(\beta_{\text{Skills}}\)         & 0.1079 & 0.1216 & 0.1223 \\
\(\beta_{\text{Certifications}}\) & 0.0721 & 0.0876 & 0.0866 \\
\(\beta_{\text{Education}}\)      & 0.2868 & 0.2778 & 0.2781 \\
\(\beta_{\text{CompanySize}}\)    & 0.4073 & 0.4011 & 0.4015 \\
\(\sigma\)                        & 0.7411 & 0.7415 & 0.7415 \\
\bottomrule
\end{tabular}
\caption{Posterior mean estimates under different prior configurations.}
\label{tab:prior_sensitivity}
\end{table}

The maximum difference across prior choices is < 0.02 for all coefficients, confirming that with N=4,000 the data overwhelms the prior. The posteriors are robust.

\begin{figure}[h]
\centering
\includegraphics[width=0.85\textwidth]{posterior-dist-prior.png}
\caption{Placeholder for prior sensitivity plots or posterior comparison visualization.}
\label{fig:prior_sensitivity_placeholder}
\end{figure}

\section{Regression Model Comparison}

The baseline linear and non-linear regression models explain roughly 44\% of the variance in salary using only 5--8 numerical predictors. In contrast, the full-feature model achieves a much higher \(R^2 = 0.9552\), showing that the inclusion of categorical variables such as job title, industry, and location substantially improves predictive performance. This result indicates that categorical factors are the dominant drivers of salary variation in this dataset.

\begin{table}[h]
\centering
\renewcommand{\arraystretch}{1.3}
\begin{tabular}{lcc}
\toprule
\textbf{Model} & \textbf{RMSE} & \textbf{\(R^2\)} \\
\midrule
\textbf{Linear (5 features)} & 27,578 & 0.4404 \\
\textbf{Non-linear (8 features)} & 27,510 & 0.4431 \\
\textbf{Full-feature (34 features)} & 7,805 & 0.9552 \\
\bottomrule
\end{tabular}
\caption{Performance comparison of Bayesian regression models.}
\label{tab:regression_comparison}
\end{table}

\begin{figure}[h]
\centering
\includegraphics[width=0.85\textwidth]{lin-reg-comp.png}
\caption{Placeholder for regression model comparison plots or posterior predictive checks.}
\label{fig:regression_comparison_placeholder}
\end{figure}

\section{Discussion}

\subsection{Issues and Observations}
\begin{itemize}
    \item The dataset appears synthetic, with uniform features and very low correlations ($< 0.03$).
    \item Numerical features give $R^2 \approx 0.44$, while categorical variables raise it to $0.96$.
    \item Classification accuracy ($\sim 52\%$) suggests remote work is nearly random.
    \item Clustering finds 4 groups, mainly separated by education and company size.
\end{itemize}

\subsection{Potential Improvements}
\begin{itemize}
    \item Use hierarchical models for group-level effects.
    \item Apply Horseshoe priors or Gaussian Processes for more flexible modeling.
    \item Extend classification to a 3-class setup (No/Hybrid/Yes).
\end{itemize}

\section{Conclusion}

The full Bayesian regression model (34 features) performed best ($R^2 = 0.9552$, RMSE = 7,805), with categorical variables explaining most salary variation. Remote work prediction was near-random ($\sim 52\%$), while clustering revealed 4 meaningful job segments. All models converged well ($\hat{R} \approx 1.0$) and showed low sensitivity to prior choice.

\section{Self-Reflection}
\begin{itemize}
    \item Learned to use PyMC's NUTS sampler while managing computation.
    \item Better understood standardization and reduced prior sensitivity in larger datasets.
    \item Gained experience evaluating models and reporting negative results clearly (classification $\approx$ random and mixed clusters).
\end{itemize}

\end{document}